\begin{helper}
Fix an activation set $A$ and a vertex $v$ in the sampling space of
\noderef{RandomIndependentSetSampling}.  Let $\nu$ be any measure satisfying
the rectangular priority clause from that definition.  If $0\le a\le 1$ and
$0\le b\le 1$, then the $\nu$-measure of the rectangle in which the activation
set is exactly $A$, the priority of $v$ is at most $b$, every activated
neighbor of $v$ has priority at most $a$, and all remaining priority
coordinates lie in $[0,1]$, is
\[
\nu[\omega_1=A]\,
b\,a^{|\{u\in V(G):u\in A\text{ and }u\sim_G v\}|}.
\]
\end{helper}
\begin{proof}
Define a threshold function $t$ on vertices by
\[
t(u)=
\begin{cases}
b, & u=v,\\
a, & u\ne v,\ u\in A,\text{ and }u\sim_G v,\\
1, & \text{otherwise}.
\end{cases}
\]
Because $0\le a\le 1$ and $0\le b\le 1$, every value of $t$ lies in
$[0,1]$.  Applying the rectangular priority clause of
\noderef{RandomIndependentSetSampling} to this threshold function gives
\[
\nu\{\omega:\omega_1=A\text{ and }0\le \omega_2(u)\le t(u)
\text{ for every }u\}
=
\nu[\omega_1=A]\prod_{u\in V(G)}t(u).
\]
The event on the left is exactly the stated rectangle: the coordinate of
$v$ is bounded above by $b$, activated neighbors of $v$ are bounded above by
$a$, and every other coordinate is bounded above by $1$.

It remains to simplify the finite product.  The factor belonging to $v$ is
$b$.  Since a simple graph has no loop at $v$, the vertex $v$ is not counted
among the vertices satisfying $u\in A$ and $u\sim_G v$.  Every such activated
neighbor contributes a factor $a$, and every remaining vertex contributes
$1$.  Hence
\[
\prod_{u\in V(G)}t(u)
=
b\,a^{|\{u\in V(G):u\in A\text{ and }u\sim_G v\}|}.
\]
Substituting this product identity into the preceding rectangular-measure
formula proves the claim.
\end{proof}
